\section{Ising Model}

\begin{exercise}[$\mathbb Z_2$ spin symmetry and the partition function]
For a finite graph $\Gamma=(V,E)$ and spins $\sigma:V\to\{\pm1\}$ define
\[
 Z_\Gamma(K,h)=\sum_\sigma
 \exp\!\left(K\sum_{ij\in E}\sigma_i\sigma_j
 +h\sum_{i\in V}\sigma_i\right).
\]
Identify the global $\mathbb Z_2$ action at $h=0$.  Expand each edge factor
into even and odd parts and prove the high-temperature loop expansion.
Expand about the two ordered configurations to obtain the dual domain-wall
expansion.  On the square lattice derive Kramers--Wannier duality
$\sinh(2K)\sinh(2K^*)=1$, including all finite-torus boundary sectors.
\end{exercise}

\begin{exercise}[Transfer-matrix decomposition]
For an $M\times N$ periodic square lattice, construct the row Hilbert space
as the tensor product of the two-dimensional spin object and define the
positive transfer operator $T_M(K)$.  Prove
$Z_{M,N}=\Tr(T_M^N)$ and that $T_M$ commutes with global spin reversal and
translation.  Introduce the Clifford algebra with the fixed convention
$\{\gamma_a,\gamma_b\}=2\delta_{ab}$, construct the Jordan--Wigner operators,
and reduce $T_M$ in each spin structure to commuting two-dimensional momentum
blocks.  Derive their rapidities
\[
 \cosh\gamma(k)=\cosh(2K^*)\cosh(2K)
 -\sinh(2K^*)\sinh(2K)\cos k.
\]
\end{exercise}

\begin{exercise}[Onsager free energy]
Put $\beta=(k_BT)^{-1}$.  Take the thermodynamic limit of the largest transfer
eigenvalue and derive
the isotropic free energy
\[
 -\beta f(K)=\frac12\log(2\sinh2K)
 +\frac1{2\pi}\int_0^\pi
 \operatorname{arcosh}(\cosh2K\coth2K-\cos\theta)\dd\theta.
\]
Locate its only nonanalytic point at
$K_c=\tfrac12\log(1+\sqrt2)$ and derive the logarithmic heat-capacity
singularity.
\end{exercise}

\begin{exercise}[Toeplitz determinants and spontaneous magnetization]
Express a long-distance spin correlation as a Toeplitz determinant.  Derive
the scalar symbol by the Clifford decomposition, factor it into inside and
outside analytic parts, and prove the strong Szeg\H{o} limit needed for its
ordered phase.  Obtain
\[
 M(K)=\lim_{r\to\infty}\langle\sigma_0\sigma_r\rangle^{1/2}
 =\left(1-\sinh^{-4}(2K)\right)^{1/8},\qquad K>K_c.
\]
\end{exercise}

\begin{exercise}[Ising characters]
Use the convention $q=e^{2\pi\ii\tau}$ from $SL_2(\Z)$.  With the theta
constants and eta function already constructed there, define
\[
 \chi_0=\frac12\left(\sqrt{\frac{\vartheta_3}{\eta}}
 +\sqrt{\frac{\vartheta_4}{\eta}}\right),\quad
 \chi_{1/2}=\frac12\left(\sqrt{\frac{\vartheta_3}{\eta}}
 -\sqrt{\frac{\vartheta_4}{\eta}}\right),\quad
 \chi_{1/16}=\frac1{\sqrt2}\sqrt{\frac{\vartheta_2}{\eta}}.
\]
Fix the branches by positive $q$-expansions.  Derive these functions from the
scaling limit of the four transfer-matrix spin structures, compute their
coefficients, and prove their $S,T$ transformation matrices.  Read off
$c=1/2$ and primary weights $0,1/2,1/16$, and construct the modular-invariant
torus partition function $\abs{\chi_0}^2+\abs{\chi_{1/2}}^2+
\abs{\chi_{1/16}}^2$.
\end{exercise}

\begin{exercise}[Ising criticality]
Restore $K=J/(k_BT)$.  Predict
\[
 k_BT_c=\frac{2J}{\log(1+\sqrt2)},\qquad
 M(T)\sim C(T_c-T)^{1/8}.
\]
On a cylinder of circumference $L$, use the three Ising characters and the
conformal Hamiltonian to derive the universal ground-energy correction
$-\pi v_s/(12L)$ and excitation gaps
$(2\pi v_s/L)(h+\bar h+n+\bar n)$.  State separately the square-lattice,
nearest-neighbor, zero-field assumptions; the measured temperature,
magnetization, and gaps; and the nonuniversal calibrations $J$ and $v_s$.
\end{exercise}

\section{Quantum Groups}

\begin{exercise}[The Hopf algebra $U_v(\mathfrak{sl}_2)$]
For $v\in\C^\times$ not a root of unity, define
$U_v(\mathfrak{sl}_2)$ by
\[
 KEK^{-1}=v^2E,\quad KFK^{-1}=v^{-2}F,\quad
 [E,F]=\frac{K-K^{-1}}{v-v^{-1}},
\]
with
\[
 \Delta K=K\otimes K,\quad
 \Delta E=E\otimes K+1\otimes E,\quad
 \Delta F=F\otimes1+K^{-1}\otimes F,
\]
and construct its counit and antipode from the Hopf identities.  Define
$[n]_v=(v^n-v^{-n})/(v-v^{-1})$ and construct every finite-dimensional
type-one irreducible $V_n$ from a highest-weight vector.  Prove the deformed
Clebsch--Gordan rule $V_m\otimes V_n\simeq
\bigoplus_{j=0}^{\min(m,n)}V_{m+n-2j}$.
\end{exercise}

\begin{exercise}[Universal $R$-matrix]
Write $K=v^H$ formally and, in the completion acting on type-one modules,
define
\[
 \mathcal R=v^{H\otimes H/2}
 \sum_{r\ge0}v^{r(r-1)/2}
 \frac{(v-v^{-1})^r}{[r]_v!}E^r\otimes F^r.
\]
Verify directly on highest-weight modules that
$\mathcal R\Delta(a)=\Delta^{\rm op}(a)\mathcal R$ and the two hexagon
identities; deduce braid-group actions and the quantum Yang--Baxter equation.
\end{exercise}

\begin{exercise}[Quantum Peter--Weyl decomposition]
Define the restricted Hopf dual
$\mathcal O_v(SL_2)$ as the span of matrix coefficients.  Prove its algebraic
Peter--Weyl decomposition
\[
 \mathcal O_v(SL_2)=\bigoplus_{n\ge0}V_n^\vee\otimes V_n
\]
and derive Schur orthogonality for the invariant positive functional when
$0<v<1$, after equipping it with the compact real-form involution and writing
the resulting Hopf $*$-algebra as $\mathcal O_v(SU(2))$.
\end{exercise}

\begin{exercise}[Basic hypergeometric and dual $q$-Hahn functions]
Put $Q=v^2$ and keep this base distinct from the modular nome.  Define
$(a;Q)_r$ and
\[
 {}_r\phi_s\!\left(\begin{matrix}a_1,\ldots,a_r\\
 b_1,\ldots,b_s\end{matrix};Q,z\right)
 =\sum_{k\ge0}\frac{(a_1,\ldots,a_r;Q)_k}
 {(Q,b_1,\ldots,b_s;Q)_k}
 \left((-1)^kQ^{k(k-1)/2}\right)^{1+s-r}z^k.
\]
Derive the $Q$-binomial and $Q$-Gauss sums by finite telescoping and analytic
continuation.  Define the dual $Q$-Hahn polynomials by
\[
 R_n(Q^{-x}+\gamma\delta Q^{x+1};\gamma,\delta,N\mid Q)
 ={}_3\phi_2\!\left(\begin{matrix}Q^{-n},Q^{-x},
 \gamma\delta Q^{x+1}\\ \gamma Q,Q^{-N}
 \end{matrix};Q,Q\right).
\]
Derive their difference equation and finite orthogonality.  Solve the
$U_v(\mathfrak{sl}_2)$ lowering recurrence and identify its normalized
Clebsch--Gordan coefficients with these polynomials after matching the three
integer labels and normalization.  Define the terminating $Q$-Racah family by
\[
 \mathcal R_n(Q^{-x}+\delta Q^{x-N})
 ={}_4\phi_3\!\left(\begin{matrix}
 Q^{-n},\alpha\beta Q^{n+1},Q^{-x},\delta Q^{x-N}\\
 \alpha Q,\beta\delta Q,Q^{-N}
 \end{matrix};Q,Q\right)
\]
and derive its finite orthogonality and recoupling identity from two
associations of a fourfold tensor product.
\end{exercise}

\begin{exercise}[Six-vertex transfer matrices and the XXZ chain]
Write $v=e^\eta$ and define, with additive spectral parameter $u$,
\[
 R(u)=\begin{pmatrix}
 \sinh(u+\eta)&0&0&0\\
 0&\sinh u&\sinh\eta&0\\
 0&\sinh\eta&\sinh u&0\\
 0&0&0&\sinh(u+\eta)
\end{pmatrix}.
\]
Evaluate the finite universal $R$-matrix on $V_1\otimes V_1$, prove that its
checked form satisfies the Hecke quadratic relation, and derive the displayed
spectral matrix by Baxterization.  Prove its Yang--Baxter equation directly
from the braid and Hecke relations.  Form the periodic row transfer matrix, prove
commutativity for different $u$, and take its logarithmic derivative at
$u=0$ to obtain, up to scalar and scale,
\[
 H_{XXZ}=\sum_{j=1}^N
 (\sigma_j^x\sigma_{j+1}^x+\sigma_j^y\sigma_{j+1}^y
 +\cosh\eta\,\sigma_j^z\sigma_{j+1}^z).
\]
Develop the algebraic Bethe ansatz and derive the Bethe equations and energy
in every fixed-magnetization sector.
\end{exercise}

\begin{exercise}[XXZ multiplet spectroscopy]
For an open $U_v(\mathfrak{sl}_2)$-symmetric spin chain, decompose three
adjacent spin sectors in the two recoupling orders.  Express the change of
channel through normalized dual $Q$-Hahn coefficients and its iterates through
$Q$-Racah polynomials.  For a weak probe transforming as $V_2$, derive the
allowed total-spin channels and their relative transition intensities as
squares of these coefficients.  For the periodic chain, use the Bethe roots
to give the corresponding resonance frequencies.  State boundary symmetry,
anisotropy $\Delta=\cosh\eta$, polarization, temperature, and line-resolution
assumptions separately from the predicted missing lines, intensity ratios,
and frequencies.
\end{exercise}

\section{Double Affine Hecke Algebras}

\begin{exercise}[Double-affine braid groups]
For the type-$A_{r-1}$ root system in
$\{x\in\R^r:\sum x_i=0\}$, construct its Weyl chambers, affine reflection
hyperplanes, and alcoves from the root datum already developed.  Identify the
affine braid group with braids in an annulus.  Show that allowing every strand
to wind around both cycles of a torus gives the double-affine braid group,
with finite braid generators $T_i$ and commuting winding families $X_j$ and
$Y_j$.  Derive its braid and winding relations directly from homotopies of
torus braids.  Impose
$(T_i-t^{1/2})(T_i+t^{-1/2})=0$ and fix the total winding commutator to $q$ to
define $\mathsf H_{q,t}(GL_r)$.  Here $q$ is a multiplicative difference
parameter, not the earlier modular nome; when this finite Hecke action is
compared with Quantum Groups use $t^{1/2}=v$, and otherwise keep $t$ and $v$
independent.
\end{exercise}

\begin{exercise}[Polynomial representation and Cherednik operators]
Let
\[
 \mathcal P_r=\C(q,t^{1/2})[x_1^{\pm1},\ldots,x_r^{\pm1}].
\]
Let $X_i$ be multiplication, let $s_i$ exchange $x_i,x_{i+1}$, define
$T_{q,x_i}f(x_1,\ldots,x_r)=f(x_1,\ldots,qx_i,\ldots,x_r)$, and put
\[
 T_i=t^{1/2}s_i+(t^{1/2}-t^{-1/2})
 \frac{s_i-1}{x_i/x_{i+1}-1},
 \quad
 (\omega f)(x)=f(qx_r,x_1,\ldots,x_{r-1}),
\]
\[
 Y_i=T_i\cdots T_{r-1}\omega T_1^{-1}\cdots T_{i-1}^{-1}.
\]
Verify preservation of Laurent polynomials, the Hecke and affine braid
relations, and $[Y_i,Y_j]=0$.  Prove their triangular action on monomials in
Bruhat-refined dominance order.  Show on symmetric Laurent polynomials that
\[
 t^{(r-1)/2}\sum_iY_i=D_{q,t},\qquad
 D_{q,t}=\sum_i\prod_{j\ne i}\frac{tx_i-x_j}{x_i-x_j}T_{q,x_i}.
\]
\end{exercise}

\begin{exercise}[Macdonald polynomials and their kernel]
For a partition $\lambda$ of length at most $r$, define $P_\lambda$ as the
monic symmetric polynomial triangular in monomial symmetric functions and
satisfying
\[
 D_{q,t}P_\lambda=\left(\sum_{i=1}^rq^{\lambda_i}t^{r-i}\right)P_\lambda.
\]
For $0<q,t<1$, prove symmetry of $D_{q,t}$ for the torus weight
\[
 \Delta_{q,t}(x)=\prod_{i\ne j}
 \frac{(x_i/x_j;q)_\infty}{(tx_i/x_j;q)_\infty}
\]
and deduce orthogonality and completeness.  In the stable symmetric-function
ring define the power-sum pairing
\[
 \langle p_\lambda,p_\mu\rangle_{q,t}
 =\delta_{\lambda\mu}z_\lambda
 \prod_j\frac{1-q^{\lambda_j}}{1-t^{\lambda_j}}.
\]
Determine the Macdonald norm from the adjoint Pieri operators, define
$Q_\lambda$ to be dual to $P_\lambda$ for this pairing, and prove the Cauchy
kernel
\[
 \sum_\lambda P_\lambda(a;q,t)Q_\lambda(b;q,t)
 =\prod_{i,j}\frac{(ta_ib_j;q)_\infty}{(a_ib_j;q)_\infty}.
\]
Construct nonsymmetric joint $Y$-eigenfunctions by intertwiners and recover
$P_\lambda$ with the finite-Hecke symmetrizer.
\end{exercise}

\begin{exercise}[Elliptic deformation and trigonometric degeneration]
For $|p|,|q|<1$ define
\[
 \theta(z;p)=(z;p)_\infty(p/z;p)_\infty,\qquad
 \Gamma(z;p,q)=\prod_{j,k\ge0}
 \frac{1-p^{j+1}q^{k+1}/z}{1-p^jq^kz}.
\]
Prove their difference and reflection identities.  Use the multiplicative
theta addition formula to construct the rank-one elliptic
difference-reflection operators and verify their dynamical braid relation.
Derive the terminating Frenkel--Turaev sum by telescoping that relation.
Take $p\to0$ to obtain the Askey--Wilson operator and its
${}_4\phi_3$ eigenpolynomials; specialize its four boundary parameters to
recover the rank-one Macdonald operator.  Then take $q=e^{-\epsilon}$,
$t=q^g$, $\epsilon\downarrow0$, to obtain the Jacobi/Jack differential
operator.  Record every parameter map so that $p$ is always the elliptic nome,
$q$ the difference step, and $t$ the Hecke parameter.
\end{exercise}

\begin{exercise}[Ruijsenaars--Schneider spectrum]
Put $x_j=e^{\ii\theta_j}$ and regard the commuting symmetric polynomials in
$Y_i$ as quantum integrals of the compact trigonometric
Ruijsenaars--Schneider system.  With the positive Macdonald weight, construct
their self-adjoint realization and prove
\[
 E_\lambda=\sum_{i=1}^rq^{\lambda_i}t^{r-i}
\]
for the first normalized Hamiltonian.  Derive the first Macdonald Pieri rule
for multiplication by $x_1+\cdots+x_r$ and hence the allowed transitions
$\lambda\to\lambda+\square$ and their squared normalized intensities.  In the
$q=e^{-\epsilon}$, $t=q^g$ limit recover the Sutherland energies and Jack
selection rule.  Thus a periodic density probe predicts exact line positions,
missing lines, and intensities; state the finite-particle, coupling, boundary,
and spectral-resolution assumptions.
\end{exercise}
